\documentclass{ximera}

\input{../../preamble.tex}

\title{A New Operation}

\begin{document}

\begin{abstract}
%%
\end{abstract}
\maketitle





With the caveat that domains need to be aligned, we have many operations on functions that make up an arithmetic on functions. 




$\blacktriangleright$ \textbf{\textcolor{blue!75!black}{Addition}}  

The sum of two functions is again a function. 

The additive identity is the zero function, $Zero(x) = 0$ for all real numbers.

For each function, $f$, there is another function $-f$, such that $f + (-f) = 0$.   $f$ and $-f$ are inverses of each other with respect to addition.




$\blacktriangleright$ \textbf{\textcolor{blue!75!black}{Multiplication}} 

The product of two functions is again a function.  

The multiplicative identity is $One(x) = 1$ for all real numbers.

For each function, $f$, there is another function $\frac{1}{f}$, such that $f \cdot \left( \frac{1}{f} \right) = 1$, again domain restrictions might be needed to avoid problems.   $f$ and $\frac{1}{f}$ are inverses of each other with respect to multiplication.  (If we were talking about numbers, then we might also use exponential notation $\frac{1}{n} = n^{-1}$.


$\blacktriangleright$ \textbf{\textcolor{blue!75!black}{Composition}} 

Composition is a new operation on functions.

The composition of two functions is again a function. 

The identity function, $Id(x) = x$ for all $x$, is the identity element with respect to composition.



\[   f \circ Id = f    \, \text{ and } \, Id \circ f = f        \]



The inverse of $f$ with respect to composition is another function whose composition with $f$ produces the identity function.

Stealing the exponential notation from numbers, the symbol for the inverse of $f$ is $f^{-1}$.

\[   f \circ f^{-1} = Id    \, \text{ and } \, f^{-1} \circ f = Id       \]



















\subsection*{Learning Outcomes}
\hfill\break

\begin{sectionOutcomes}
In this section, students will 

\begin{itemize}
\item use function notation.
\item evaluate functions.
\item solve equations involving functions.
\item analyze functions.
\item focus on numeric functions.
\end{itemize}
\end{sectionOutcomes}









\pdfOnly{
\begin{center}
ooooo-=-=-=-ooOoo-=-=-=-ooooo

more examples can be found in the online edition.
\end{center}
}




\begin{onlineOnly}
\begin{center}
\textbf{\textcolor{green!50!black}{ooooo-=-=-=-ooOoo-=-=-=-ooooo}} \\

more examples can be found by following this link\\ \link[More Examples of Composition]{https://ximera.osu.edu/csccmathematics/precalculus/precalculus/composition/examples/exampleList}

\end{center}
\end{onlineOnly}


















\end{document}
